\documentclass{article}
\usepackage{graphicx} % Required for inserting images
\usepackage{James}
\usepackage{dynkin-diagrams}
\title{Category Theory}
\author{James Hindmarch}
\setcounter{section}{0}
\date{October 2023}
\begin{document}

\maketitle
\tableofcontents

\newpage
\newlec


%LECTURE 1




\section{Chapter 1: Definitions \& Examples}
\begin{fancydefine}{Category}
    A category $\calC$ consists of:
    \begin{itemize}
        \item A collection of objects, denoted $\ob \calC$. ($A$, $B$, $C$, $\dots$)
        \item A collection of morphisms, $\mor \calC$. ($f$, $g$, $h$, $\dots$)
        \item Two operations $\dom , \, \cod : \mor \calC \to \ob \calC$. We write $f: A \to B$ or $A \xrightarrow[]{f} B$ to mean $\dom f = A$, $\cod f = B$.
        \item An operation $A \mapsto (1_{A} : A \to A)$ (i.e. an identity associated with each object).
        \item A partial binary operation $(f,g) \mapsto fg$ called composition if and only if $\dom f = \cod g$, and then $\dom(fg)=\dom(g)$, $\cod(fg)=\cod(f)$. This satisfies: \begin{itemize}
            \item $f1_{A} = f$, $1_{A}g=g$ whenever such composites are defined.
            \item $f(gh) = (fg)h$ whenever $fg$ and $gh$ are defined.
        \end{itemize}
    \end{itemize}
\end{fancydefine}
\remarks \hspace{0.1cm}
\begin{enumerate}
    \item[(1)] We don't require $\ob \calC$ to be sets, nor $\mor \calC$ (but if they are, we call $\calC$ a \textit{small category}).
    \item[(2)] Objects are not technically necessary at all in this definition, since we could just say there are a bunch of $\id$ morphisms (but this would be counterintuitive and we won't do it)
    \item[(3)] $fg$ means, as always, first $g$, then $f$. Computer scientists haven't understood this yet, so some books write this differently.
\end{enumerate}
\exs \hspace{0.1cm}
\begin{enumerate}
    \item[(a)] $\cat{Set} = $ category of all sets, and functions between them. (Strictly, set theory does not specify a codomain (although it should), so we \textit{technically} require a morphism to be $(f,B)$ defining a function and a codomain, but usually it's clear from context).
    \item[(b)] \begin{itemize}
        \item $\cat{Grp} = $ category of groups and group homomorphisms
        \item $\cat{Ring} = $ category of rings and ring homomorphisms
        \item $\cat{Vect$_{k}$} = $ category of $k$-vector-spaces and $k$-linear maps
    \end{itemize}
    \item[(c)] \begin{itemize}
        \item $\cat{Top} = $ category of topological spaces and continuous maps between them
        \item $\cat{Met} = $ category of metric spaces and (in this course) non-expansive maps between them (so that isomorphisms are isometries).
        \item $\cat{Mfd} = $ category of smooth maps and $C^{\infty}$ maps
        \item $\cat{TopGrp} = $ category of topological groups and continuous group homomorphisms (we can combine categories like this in other, similarly obvious ways, as long as it makes sense).
    \end{itemize}
    \item[(d)] $\cat{Htpy} = $ category of topological spaces, and morphisms given by \textit{homotopy classes of} continuous maps.\gap More generally, if $\sim$ is an equivalence relation on $\mor \calC$ such that \[ f \sim g \Rightarrow \dom f = \dom g, \, \cod f = \cod g \] and such that \[ fh \sim gh, \, kf \sim kg \] whenever defined, then we call $\sim$ a \textit{congruence} on $\calC$, and have a \textit{quotient category}, $\frac{\calC}{\sim}$ with: \[
    \ob \\calC/{\sim} = \ob \calC
    \]
    \[ \mor (\calC/{\sim}) = (\mor \calC)/{\sim} \]
    \item[(e)] $\cat{Rel} = $ category of sets and \textit{relations} $\calR \subseteq A \times B$ (where $A$ and $B$ are specified) with composition being defined by \[ \calS \circ \calR = \{ (a,c) \, | \, (\exists b \in B) \, \, (a,b) \in \calR, \, (b,c) \in \calS \} \] Note that \cat{Set} is a subcategory of \cat{Rel}. There is also $\cat{Part} = $ category of all sets and partial functions (functions from subsets of $A$ to $B$), and so we have \[
    \cat{Set} \subset \cat{Part} \subset \cat{Rel}. \]
    \item[(f)] Given any category $\calC$, we define $\calC^{\text{op}}$ (the \textit{opposite} of $\calC$) to have the same objects and morphisms, but $\dom$ and $\cod$ are interchanged, and composition is reversed.
    \item[(g)] A small category with \underline{one} object is called a \textit{monoid} (i.e. a \textit{semigroup} with a 1). In particular, a group is a small category with one object in which every morphism is an isomorphism.
    \item[(h)] A \textit{groupoid} is a category in which every morphism is an isomorphism. For example, the \underline{fundamental groupoid} $\pi_{1}(X)$ of a space has \[ \ob \pi_{1}(X) = X \]
    and morphisms $x \to y$ are homotopy classes of paths from $x$ to $y$.
    \item[(i)] A category $\calC$ with \underline{at most} one morphism $A \to B$ between any pair of objects is a \textit{preorder} (i.e. a class of objects with a reflexive, transitive relation). \gap In particular, a Poset (partially ordered set) is a small category which is a preorder in which the only isomorphisms are the identity morphisms. \gap "It is nice to use posets as an example whenever you are presented with a theorem about categories which you do not understand, since they are familiar objects, and theorems are usually not entirely trivial on posets."
    \item[(j)] Given $k$ a field, have $\cat{Mat$_{k}$} = $ The category with $\ob \cat{Mat$_{k}$} = \bbN$, and morphisms $n \to p$ are $p\times n$ matrices in $k$, composition is obviously matrix multiplication. Identities are $\id_{n \times n}$.
\end{enumerate}





\newlec


%LECTURE 2









\begin{fancydefine}{Functor}
Let $\calC$ and $\calD$ be categories, a \textbf{functor} \[ F : \calC \to \calD\]
consists of mappings \begin{align*}
    \ob \calC &\xrightarrow[]{F} \ob \calD\\
    \mor \calC &\xrightarrow[]{F} \mor \calD\\
\end{align*}
satisfying
\begin{align*}
    F(\dom f) &= \dom (Ff)\\
    F(\cod f) &= \cod (Ff)\\
    F(1_{A}) &= 1_{FA}\\
    F(fg) &= (Ff)(Fg).
\end{align*}
\end{fancydefine}
\exs \textbf{Examples of Functors}
\begin{enumerate}
    \item[(a)] Forgetful Functors - These are functors which 'forget' some structure; there's no formal definition, but you know them when you see them.
    \begin{itemize}
        \item $\cat{Grp} \to \cat{Set}$
        \item $\cat{Rng} \to \cat{Set}$
        \item $\cat{Top} \to \cat{Set}$
        \item $\cat{Rng} \to \cat{AbGrp}$
        \item $\cat{Met} \to \cat{Top}$
        \item $\cat{TopGrp} \to \cat{Top}$
        \item $\cat{TopGrp} \to \cat{Grp}$
    \end{itemize}
    \item[(b)] The free group construction \[ A \mapsto FA \] such that any mapping \[ f : \to UG \] to the underlying set of a group extends uniquely to a homomorphism \[ FA \to G.\] We make this into a functor \[ F: \cat{Set} \to \cat{Grp}, \] and given \[ f : A \to B, \] $Ff$ is the unique homomorphism extending \[ A \xrightarrow[]{f}  B \to FB. \] Given $g: B \to C$, $F(gf)$ and $F(g)F(f)$ both extend the same mapping $A \to FC$, so they're equal.
    \item[(c)] The powerset construction defines a functor \[ P : \cat{Set} \to \cat{Set}, \] where $PA = \{ \text{Subsets of } A \}$, and if $f: A \to B$, \[Pf(A') = \{f(a) \, | \, a \in A'\}.\]
    \item[(d)] We also have a \textit{totally unrelated} functor \[P^{*} : \cat{Set}^{\text{op}} \to \cat{Set},\] (or $\cat{Set} \to \cat{Set}^{\text{op}}$), which is given by $P^{*}A = PA$ and for $f : A \to B$, \[ P^{*}f(B') = \{a \in A \, | \, f(a) \in B'\}. \] We call $P^{*}$ a \textbf{contravariant} functor - this means that it 'reverses the arrows' in the category (i.e. it goes from a category to its opposite category), the opposite of a contravariant functor is a \textbf{covariant} functor.
    \item[(e)] The construction of dual spaces is a functor: \[ (-)^{*}: \cat{Vect$_{k}$}^{\text{op}} \to \cat{Vect$_{k}$},\] $V^{*}$ is the space of linear mappins \[ V \to k \], if $f: V \to W$, then $f^{*}: W^{*} \to V^{*}$ is given by composition with $f$.
    \item[(f)] The assignment $\calC \to \calC^{\text{op}}$ defines a functor \[ \cat{Cat} \to \cat{Cat}, \] where $\cat{Cat}$ is the category of \textit{small} categories (to avoid any Russell's paradox problems) and functors between them.
    \item[(g)] A functor between monoids is monoid homomorphism. \gap A covariant functor between posets is an order-preserving map (contravariant: order-reversing). \gap Let G be a group (considered as a $1$-object category), then a functor \[ F : G \to \cat{Set} \] consists of a set $A (=F(*))$ together with mappings \begin{align*}
        A &\to A\\
        a &\mapsto g\cdot a
    \end{align*}
    for the morphisms, which are compatible with identity and composition (i.e. it is an \underline{action} or \underline{permutation representation} of $G$ on the set $A$). \gap In a similar way, a functor $G \rightarrow \cat{Vect$_{k}$}$ is a $k$-linear representation of $G$
    \item[(h)] The fundamental group defines a functor \[ \pi_{1} : \cat{Top}^{*} \to \cat{Grp},\] where $\cat{Top}^{*}$ is the category of topological with basepoint, and basepoint preserving maps. We could similarly have a functor from the category \textit{un}pointed topological spaces $\cat{Top}$ to the category $\cat{Groupoid}$.
\end{enumerate}
\begin{fancydefine}{Natural Transformation}
    Let $\calC$ and $\calD$ be categories, $F,G: \calC \rightrightarrows \calD$ two functors, then a \textbf{natural transformation} is a mapping \begin{align*}
        \alpha : F \to G
    \end{align*}
    which takes
    \begin{align*}
        \ob \calC \to \mor \calD
    \end{align*}
    which we write $A \to \alpha_{A}$. Moreover, this mapping satisfies \[ \alpha_{A}: FA \to GA \] for all $A$, and for any $f: A \to B$ in $\calC$, the following square commutes
    \begin{center}
        \begin{tikzcd}
            FA \arrow{d}{\alpha_{A}} \arrow{r}{Ff}
            & FB \arrow{d}{\alpha_{B}} \\
            GA \arrow{r}{Gf}
            & GB
            \end{tikzcd}
\end{center}
\end{fancydefine}
If we also have $\beta : G \to H$ we can then define $(\beta \alpha)$ by \[ (\beta \alpha)_{A} = \beta_{A}\alpha_{A}.\]
So we have a category $[\calC, \calD]$ whose objects are functors $\calC$ to $\calD$ and whose morphisms are natural transformations.
\exs
\begin{enumerate}
    \item[(a)] Given a vector space $V$, we have a linear map \[ \alpha_{V} : V \to V^{**}\] sending $v \in V$ to \[ \text{evaluate at $v$} : V^{*} \to k\] This is a natural transformation \[ 1_{\cat{Vect$_{k}$}} \to (-)^{**}.\]
    \item[(b)] We have a mapping of $\cat{Set}$, $\alpha_{A} : A \to PA$ which sends $a \in A$ to $\{ a \}$. This is a natural transformation \[ 1_{\cat{Set}} \to P.\] since $Pf(\{a\}) = \{f(a)\}$.
    \item[(c)] The mapping sending $A$ to the inclusion of generators $A \to FA$ is a natural transformation from \[ 1_{\cat{Set}} \to UF\] where $F$ is the free group functor and $U$ is the forgetful functor.
    \item[(d)] If $f,g : P \rightrightarrows Q$ are order-preserving maps of posets, there exists a (unique) natural transformation \[ f \to g\] if and only if $f(p) \leq g(p) \, \forall p$.
    \item[(e)] Suppose we are given two group homomorphisms \[ u,v: G \rightrightarrows H. \] A natural transformation $u \to v$ is an element $h \in H$ such that \[ h\cdot u(g) = v(g) \cdot h \, \, \forall g \in G, \] or equivalently: \[ v(g) = h \cdot u(g) \cdot h^{-1}.\] That is to say, natural transformations are exactly conjugacy relations. In particular, the natural transformations $u \to v$ are the elements of the centralizer $u(G)$.
    \item[(f)] If $A$ and $B$ are permutation representations of $G$, then a natural transformation \[ A \xrightarrow{f} B\] is a mapping satisfying \[ g \cdot f(a) = f(g \cdot a)\] for all $a \in A, \, g \in G$.
    \item[(h)] For a (nice) space $X$ with basepoint $x$, there is this thing that we should mention, even though it won't come up again (unless you're taking algebraic topology),called the Hurewicz Homomorphism \[ \cat{Top$^{*}$} \to \cat{Top}, \] given by \[ h_{n,x} : \pi_{n}(X,x) \to H_{n}(X).\] This is a natural transformation from $\pi_{n}$ to $H_{n}U$ where $U$ forgets the basepoint
\end{enumerate}







%LECTURE 3





\newlec



\begin{fancydefine}{Isomorphism of Categories}
    We say that two categories $\calC$ and $\calD$ are \textbf{isomorphic} if there exist functors \[ F: \calC \to \calD \] \[ G: \calD \to \calC\] such that $GF=1_{\calC}$ and $FG=1_{\calD}$.\gap For example $\cat{Rel} \cong \cat{Rel}^{\text{op}}$ via the functor: \begin{align*}
        A &\mapsto A\\
        R &\mapsto R^{\text{op}} = \{(b,a) \, : \, (a,b) \in R\}
    \end{align*}
\end{fancydefine}
However, this notion of equivalence of categories turns out to be not super useful, and there's a weaker notion of equivalence that is used more frequently. To define it, we need the following:
\lem Let $\alpha : F \to G$ be a natural transformation between two functors \[ F, G : \calC \rightrightarrows \calD. \] Then $\alpha$ is an isomorphism in $[\calC, \calD] \Leftrightarrow $ each $\alpha_{A}$ is an isomorphism in $\calD$.
\proof \hspace{0.1cm} \\
$\Rightarrow$: Obvious, since composition in $[\calC, \calD]$ is pointwise \gap 
$\Leftarrow$: Suppose $\beta_{A}$ is inverse to $\alpha_{A}$ for each $A$, then Given $f: A \to B$ in $\calC$, consider:

\begin{center}
    \begin{tikzcd}
        FA \arrow[d,xshift=-0.7ex,"\alpha_{A}",swap] \arrow{r}{Ff}
        & FB \arrow[d,xshift=-0.7ex,"\alpha_{B}",swap] \\
        GA \arrow{r}{Gf} \arrow[u,xshift=0.7ex,"\beta_{A}",swap]
        & GB \arrow[u,xshift=0.7ex,"\beta_{B}",swap]
        \end{tikzcd}
\end{center}
Then $(Ff)\beta_{A} = \beta_{B} \circ \alpha_{B} (Ff)\beta_{A} = \beta_{B}(Gf)\alpha_{A} \circ \beta_{A}=\beta_{B}(Gf).$ So we have that $\beta$ is a natural transformation, and hence an inverse for $\alpha$ in $[\calC, \calD]$. \endproof
\begin{fancydefine}{Equivalence of categories}
    Let $\calC$ and $\calD$ be categories, by an \textbf{equivalence} between $\calC$ and $\calD$, we mean:
    \begin{itemize}
        \item A pair of functors $F: \calC \to \calD$, $G : \calD \to \calC$.
        \item Natural isomorphisms: \begin{align*}
            \alpha : 1_{\calC} \to GF & & \beta : FG \to 1_{\calD}
        \end{align*}
        (We've switched the directions around, but don't worry, we will see why in Chapter 3)
    \end{itemize}
    We write $\calC \simeq \calD$ if there exists an equivalence between $\calC$ and $\calD$.
    
\end{fancydefine}

\textbf{Definition. }(Categorical Property) We say a property $P$ of categories is \textbf{Categorical} if it is preserved under equivalence. (If $\calC$ has $P$, and $\calC \simeq \calD$ then $\calD$ has $P$.)
\begin{itemize}
    \item Being a preorder \underline{is}
    \item Being a partial order \underline{is not}
    \item Being a groupoid \underline{is} (I think)
    \item Being a group \underline{is not}
\end{itemize}
\exs \hspace{0.1cm} \begin{enumerate}
    \item[(a)] Let \cat{Set$^{*}$} be the category of sets with base point and basepoint preserving functions. Then \[\cat{Set$^{*}$} \simeq \cat{Part}\] via \[ F: \cat{Set$^{*}$} \to \cat{Part}\] defined by \[ F(A,a)= A \setminus \{a\}\] \[ F((A,a) \xrightarrow f (B,b))(x) = \begin{cases}
        f(x), & f(x) \neq b\\
        \text{undefined}, & \text{o/w}
    \end{cases}\]
    and via \[ G : \cat{Part} \to \cat{Set$^{*}$} \] defined by:
    \[ G(A) = (A \cup \{ A \}, A),\] \[ G(A \rightharpoondown B)(x) = \begin{cases}
        f(x), & \text{if $x \in A$ and $f$ defined}\\
        B, & \text{o/w}.
    \end{cases}
    \] ($\rightharpoondown$ means partial function.) Then $FG = 1_{\cat{Part}}$, since what we're doing is: adding a basepoint, then throwing away that basepoint; changing the function, then throwing away those changes. $GF \neq 1_{\cat{Set$^{*}$}}$, however there is an obvious isomorphism (apparently, need to look this up). 
    There can be no isomorphism between these categories, however, since isomorphisms preserve isomorphism classes within categories, and we have $\emptyset$ in \cat{Part} a $1$-object isomorphism class, and no such $1$-object isomorphism class in $\cat{Set$^{*}$}$
    \item[(b)] Let \cat{fdVect$_{k}$} be the category of finite dimensuonal vector spaces over $k$. \gap Then \[\cat{fdVect$_{k}$} \simeq \cat{fdVect$_{k}^{\text{op}}$} \] via \[ F = (-)^{*} : \cat{fdVect$_{k}$} \to \cat{fdVect$_{k}^{\text{op}}$}\] and \[ G = (-)^{*} : \cat{fdVect\subk}^{*} \to \cat{fdVect\subk\opp}. \] $\alpha$ and $\beta$ are both the natural transformation of $1.7(a)$.
    \item[(c)] $\cat{fdVect\subk} \simeq \cat{Mat$_{k}$}$ of $1.3(j)$. via \[ F : \cat{Mat\subk} \to \cat{fdVect\subk} \] which sends $n$ to $k^{n}$ and the matrix to the linear map $k^{n} \to k^{p}$ given by the matrix, and \[ G: \cat{fdVect\subk} \to \cat{Mat\subk} \] where we have to choose a basis for each finite dimensional vector space. 
    Then $G(V) = \dim V$ and $G(V \xrightarrow f W) = \text{ Matrix of $f$ in a chosen basis}$. Then $GF=1_{\cat{Mat\subk}}$ provided we make the choice of the standard basis, and $FG \simeq 1_{\cat{fdVect\subk}}$ since chosen bases define an isomorphism: \[ k^{\dim V} \to V \] which are natural on $V$.
\end{enumerate}
\begin{fancydefine}{Properties of Functors}
    Let \[ F : \calC \to \calD\] be a functor.
    \begin{enumerate}
        \item[(a)] We say $F$ is \textbf{faithful} if, given $f,g \in \mor \calC$, we have \begin{align*}
            \dom f = \dom g, & & \cod f = \cod g, & & \text{and} \, \, \,   Ff=Fg \, \Rightarrow \, f=g.
        \end{align*}
        \item[(b)] We say $F$ is \textbf{full} if, given \[ FA \xrightarrow g FB, \] there exists $f : A \to B$ with $Ff = g$
        \item[(c)] We say $F$ is \textbf{essentially surjective} if every $B \in \ob \calD$ is isomorphic to some $FA$.
    \end{enumerate}
    Note:  Essential injectivity is implied by full and faithful. That is, if $FA \xrightarrow g FB$ is an isomorphism, then the unique $f : A \to B$ with $Ff = g$ is an isomorphism. (With inverse the unique $B \to A$ which gets mapped to $g^{-1}$ by $F$.)
\end{fancydefine}
\lem Let \[ F: \calC \to \calD \] be a functor. Then $F$ is part of an equivalence $calC \simeq \calD \, \Leftrightarrow$ $F$ is full, faithful, and essnetially surjective.
\proof $\Rightarrow$: Suppose, given $G,\alpha, \beta$ as in 1.9, the existence of $\beta$ ensures that \[ B \simeq FGB \] for any $B \in \ob \calD$ (so we have essential surjectivity). Also, we have for any $A \xrightarrow f B$ in $\calC$: \[f = \alpha_{B}^{-1}(GFf)\alpha_{A}, \] so $f$ is determined by $\dom f, \cod f,$ and $Ff$; so we have faithfulness. Then for fullness, given \[ FA \xrightarrow g FB, \] we define \[ f = \alpha_{B}^{-1}(Gg)\alpha_{A} : A \to B, \] then \[ GFf = Gg. \] But $G$ is faithful, since $F$ is (and the same arguments apply), so $Ff=g$.\gap 
$\Leftarrow$: For each $B \in \ob \calD$, choose an isomorphism \[ \beta_{B} : FGB \to B. \] This defines $G$ on objects. To define $G$ on morphisms, we take $G(B \xrightarrow g C)$ to be the unique $GB \to GC$ whose image under $F$ is \[ FGB \xrightarrow \beta_{B} B \xrightarrow g C \xrightarrow \beta_{C}^{-1} FGC.\] This is functorial since given $h: C \to D$ \[ G(hg) \text{   and   } (Gh)(Gg)\] have the same image under $F$, so they must be equal by faithfulness. 
Moreover $\beta_{B}$ is a natural transformation, by construction: \[ FG \to 1_{\calD}.\] So just need to construct $\alpha : 1_{\calC} \to GF$, we do this by: \[ \alpha_{A} : A \to GFA \] to be the unique morphism, whose image under $F$ is \[ FA \xrightarrow \beta_{FA}^{-1} FGFA. \] (This is an isomorphism since the image under $F$ is an isomorphism). Then, the naturality squares for $\alpha$ are mapped by $F$ to naturality squares for $\beta^{-1}$. So they commute. \endproof
\define[Skeletal Category] We say $\calC$ is \textbf{skeletal} if every isomorphism class in $\ob \calC$ has just one member. By a \textbf{skeleton} of $\calC$, we mean a full subcategory $\calC'$ containing just one object from each isomorphism class. For example \cat{Mat\subk} is skeletal. Note that any equivalence between two skeletal categories is (bijective on objects, and hence) an isomorphism. \gap 









%Lecture 4
\newlec 




Skeletons are basically a diversion, and are fundamentally unimportant. \gap We want something \textit{like} injectivity and surjectivity for morphisms within a category, but can't use those words, because the objects need not be sets. So:
\begin{fancydefine}{Monomorphism/Epimorphism/Balanced Category}
    Let $f: A \to B$ be a morphism in $\calC$. We say: \begin{itemize}
        \item $f$ is a \textbf{monomorphism} (or \textbf{monic}) if \[fg = fh \Rightarrow g=h,\] whenever composition is defined.
        \item Dually, $f$ is an \textbf{epimorphism} or (\large\textbf{EPIC}\normalsize) if \[ gf=hf \Rightarrow g=h,\] again, whenever composition is defined.
    \end{itemize}
    If $f$ is an isomorphism then it must be \textbf{monic} and \textbf{epic}. We say $\calC$ is a \textbf{balanced} category if the converse holds. \gap $f : A \rightarrowtail B$ means $f$ is monic; $A \twoheadrightarrow B$ means $f$ is epic.  
    
\end{fancydefine}
\exs \begin{enumerate}
    \item[(a)] In \cat{Set}, monic $\Leftrightarrow$ injective ($\Leftarrow$ obvious, for $\Rightarrow$, use morphisms $1 \to A$).\\ Epic $\Leftrightarrow$ surjective ( $Leftarrow$ easy, for $\Rightarrow$, use morphisms $B \to 2 = \{0,1\})$. So $\cat{Set}$ is balanced.
    \item[(b)] In $\cat{Gp}$, monic $\Leftrightarrow$ injective (Use $\bbZ$ in place of $1$).\\ And epic $\Leftrightarrow$ surjective (proof nontrivial). So $\cat{Grp}$ is balanced.
    \item[(c)] In $\cat{Rng}$, monic $\Leftrightarrow$ injective, but $\bbZ \hookrightarrow \bbQ$ is epic, so $\cat{Rng}$ isn't balanced.
    \item[(d)] In \cat{Top}, monic $\Leftrightarrow$ injective, and epic $\Leftrightarrow$ surjective, but \cat{Top} isn't balanced.
    \item[(e)] In a preorder, all morphisms are monic and epic, but a preorder is balanced iff it's an equivalence relation. 
\end{enumerate}
\section{Chapter 2: The Yoneda Lemma}
\begin{fancydefine}{Locally Small}
    We say a category $\calC$ is \textbf{locally small} if, for all $A, B \in \ob \calC$, the morphisms $A \to B$ in $\calC$ are parameterised by a set $\calC(A,B)$. Sometimes this is taken to be part of the definition of a category, but we won't do that, since P.T. Johnstone doesn't like it, and wants category theory's definition to be independent of set theory (and other reasons).
\end{fancydefine}
Given an object $A$ of a locally small category $\calC$, we have a functor \[ \calC(A, - ) : \calC \to \cat{Set}\] sending $B$ to $\calC(A,B)$ and $B \xrightarrow f C$ to the mapping $g \mapsto fg$ where $g \in \calC(A,B)$. This is functorial since composition in $\calC$ is associative. Similarly, for fixed $B$, we have a functor \[ \calC( - , B): \calC^{\text{op}} \to \cat{Set}.\] 
\begin{fancylem}{Yoneda}
     Let $\calC$ be locally small, let $A \in \ob \calC$, and $F: \calC \to \cat{Set}$ be a functor. Then \begin{enumerate}
    \item[(i)] There is a bijection between natural transformations \[ \calC(A, - ) \to F \] and elements of $FA$
    \item[(ii)] Moreover, this bijection is natural in both $A$ and $F$. 
\end{enumerate}
\end{fancylem}
\proof \hspace{0.1cm} \begin{enumerate}
    \item[(i)] Given $\alpha : \calC(A, - ) \to F$, we define \[\Phi(\alpha) = \alpha_{A}(1_{A}) \in FA. \] Given $x \in FA$, we define $\Psi(x) : \calC(A, - ) \to F$ by \[ \Psi(x)_{B}(A \xrightarrow f B) = (Ff)(x) \in FB.\] This is natural, since $F$ is a functor.\gap Clearly $\Phi \Psi (x) = \Psi(x)_{A}(1_{A}) = F(1_{A})(x) = x$.\\ If $\alpha: \calC(A, - ) \to F$, then \[\Psi \Phi(\alpha)_{B}(f) = \Psi(\alpha_{A}(1_{A}))_{B}(f) = (Ff)(\alpha_{A}(1_{A})) = \alpha_{B}(f 1_{A}) = \alpha_{B}(f).\]
\end{enumerate}
\endproof
\corollary[Yoneda Embedding] For any locally small $\calC$, the mapping $A \mapsto \calC(A, - )$ is a full and faithful functor \[ \calC^{\text{op}} \to [\calC, \cat{Set}].\]
\proof Putting $F=\calC(B, - )$ in 2.2, we get a bijection from $\calC(B,A)$ to the set of natural transformations \[ \calC(A, - ) \to \calC(B, -).\] This sends $f: B \to A$ to the natural transformation given by composition with $f$. Functorial since composition in $\calC$ is associative. \endproof 
This says that any locally small category is equivalent toa full subcategory of a functor category $[\calC, \cat{Set}]$. This is like Cayley's theorem in Group Theory (which is technically a special case of this result), which takes this abstract thing called a group, and says we can present it as an action on a set of objects. Here we are saying we can present categories within the catgegory of sets (so long as the category is locally small).
\begin{fancydefine}{Representable/Representation of a functor}
    If $\calC$ is locally small, a functor \[ F: \calC \to \cat{Set}\] is called \textbf{representable} if it is isomorphic to $\calC(A, -)$ for some $A \in \ob \calC$. \gap By a \textbf{representation} of $F$, we mean a pair $(A,x)$ where $A \in \ob \calC$. and $x \in FA$ is such that $\Psi(x)$ is a natural isomorphism $\calC(A,-) \to F$. ($x$ is called a \textbf{universal element} of $F$).
\end{fancydefine}
\corollary If $(A,x)$ and $(B,y)$ are both representations of a functor $F$, then there is a unique isomorphism $f : A \to B$ which is compatible with $F$, i.e. such that \[ Ff(x)=y.\]
\proof $(Ff)(x)=y$ is equivalent to saying that 
\begin{center}

    \begin{tikzcd}[column sep=scriptsize]
        \calC(B, -) \arrow[dr, "\Phi(y)", swap] \arrow[rr, "\calC (f{,} -)"] & & \calC(A, -) \arrow[dl, "\Psi(x)"] \\
        & F
    \end{tikzcd}
\end{center}
commutes. So this holds if and only if $f$ is the unique ismorphism sent by the Yoneda embedding to $\Psi(x)^{-1}\Phi(y)$. \endproof
\proof 2.2(ii) Suppose for a moment that $\calC$ is small, so that $[\calC, \cat{Set}]$ is locally small, then we have two functors \[ \calC \times [\calC, \cat{Set}] \to \cat{Set}.\] One of them sends $(A,F)$ to $FA$, the other is the composition: \[ \calC \times [\calC, \cat{Set}] \xrightarrow{Y \times 1} [\calC, \cat{Set}]^{\text{op}} \times [\calC, \cat{Set}] \xrightarrow{[\calC,\cat{Set}](-,-)} \cat{Set}.\] where $Y$ is the Yoneda embedding of $\calC$.
In this case, Yoneda (ii) says that $\Phi$ and $\Psi$ are natural transformations between these functors. \gap Now, once we write it out, we realize we don't \textit{need} $\calC$ to be locally small, in elementary terms, we're simply saying that Given $f: A \to A'$ and $\alpha: F \to F'$, and $x \in FA$, if $x'=\alpha_{A}(Ff)(x)=(F'f)\alpha_{A}(x)$, then $\Psi(x')$ is the composite \[ \calC(A', -) \xrightarrow{\calC(f,-)} \calC(A,-) \xrightarrow{\Psi(x)} F \xrightarrow{\alpha} F. \] But this composite sends \[1_{A'} \mapsto f \mapsto (Ff)(x) \mapsto \alpha_{A'}(Ff)(x)\].
\begin{ex} \newlec 
    \begin{enumerate}
        \item[(a)] The forgetful functor $\cat{Gp} \to \cat{Set}$ is representable by $\bbZ$.\\ Similar $\cat{Rng} \to \cat{Set}$ is representable by $\bbZ[X]$.\\ $\cat{Top} \to \cat{Set}$ is representable by the $1$-point space.
        \item[(b)] $P^{*} : \cat{Set}^{\text{op}} \to \cat{Set}$ is representable by $2=\{0,1\}$ via the bijection which sends $f:A \to 2$ to $f^{-1}(1)$.\gap But $P : \cat{Set} \to \cat{Set}$ is not representable, we have $\cat{Set}(A,1) \equiv 1$ for any $A$, but $P1 \equiv 2$.
        \item[(c)] We have a functor $\Omega : \cat{Top}^{\text{op}} \to \cat{Set}$ sending a space to its set of open subsets, and a continuous map $f$ to the effect of $f^{-1}$ on open subsets. This is representable by the \underline{Sierpinski Space} $\Sigma = \{0,1\}$ with $\{1\}$ open since continuous maps $X \to \Sigma$ are exactly the characteristic functions of open subsets of $X$.
        \item[(d)] The dual vector space functor can't be representable, since it doesn't map to sets, but if we compose with the forgetful functor, \[ \cat{Vect}_{k}^{\text{op}} \xrightarrow{(-)^{*}} \cat{Vect}{k} \xrightarrow{U} \cat{Set}\] is representable by the 1-dimensional space $k$.
        \item[(e)] If $G$ is a group, the unique representable functor $G \to \cat{Set}$ is the \underline{Cayley representation}, i.e. the set $G$ with $G$ acting by multiplication.
        \item[(f)] Suppose, given two objects $A$ and $B$ of a locally small category $\calC$ and a functor $\calC^{\text{op}} \to \cat{Set}$ which sends $C$ to the cartesian product \[ C \mapsto \calC(C,A) \times \calC(C,B).\] If this is representable, we call the represented object a (categorical) \textbf{product} of $A$ and $B$, and we denote it $A \times B$\footnote{Which is notation to be careful about, since it \textit{doesn't} mean the cartesian product, although in many categories that is how it manifests. But you need to be careful because sometimes it just is not the cartesian product.}, with universal element $(pi_{1}: A \times B \to A, \,\, \pi_{2} : A \times B \to B)$, the \textbf{product projections}. This has the property: given any pair $(f:C \to A, \, g: C \to B)$, there exists a unique $h=(f,g) : C \to A \times B$ satisfying $pi_{1}h = f, \, \pi_{2}h=g.$ \gap Dually, we have the notion of \textbf{coproduct} $A + B$, with \textbf{coprojections} $\nu_{1} : A \to A+B, \, \nu_{2} : B \to A+B.$
        \item[(g)] Given a parallel pair $f,g: A \rightrightarrows B$ in a locally small category $\calC$, we have a functor $F: \calC^{\text{op}} \to \cat{Set}$ which sends $C$ to $\{h: C \to A | fh = gh\}$ A representation of this functor is called an \textbf{equalizer} for $(f,g)$: it consists of an object $E$ with $e : E \to A$ satisfying $fe = ge$ and such that any $h$ with $fh=gh$ factors uniquely through $e$. Note that $e$ is necessarily a monomorphism, we say a monomorphism is \textbf{regular} if it occurs as an equalizer.
    \end{enumerate}
    In $\cat{Set}$, products are cartesian products, and coproducts are disjoint unions $A \times \{0\} \cup B \times \{ 0 \}$ \gap In $\cat{Gp}$, products are cartesian products, but the coproduct $A+B$ is the \underline{free product} $A * B$.\gap In $\cat{Set}$ or $\cat{Gp}$, the equalizer of $f,g : A \rightrightarrows B$ is the inclusion $\{a \in A | f(a)=g(a)\} \hookrightarrow A$. \gap The coequalizer of $(f,g)$ is the quotient $\frac{B}{R}$ where $R$ is the smallest \[\begin{cases} \text{Equivalence Relation} & \text{in \cat{Set}}\\
        \text{Congruence} & \text{in \cat{Gp}} \end{cases}\]
    which contains all pairs $(f(a),g(a)), a \in A$. \gap In both these categories, all monomorphisms and all epimorphisms are regular, but in $\cat{Top}$ not all injections/surjections are regular monomorphisms/epimorphisms.
\end{ex}
\begin{fancydefine}{Separating/Detecting Families}
    Let $\calC$ be a locally small category. Let $\calG$ be a class of objects of $\calC$, then \begin{enumerate}
        \item[(a)] We say $\calG$ is a \textbf{separating family} for $\calC$ if the functors $\calC(G,-)$, $G \in \calG$ are collectively faithfulm i.e. if given $f,g : \rightrightarrows B$, the equations $fh=gh$ fo all $h:G \to A$ with $G \in \calG$ implies that $f=g$.
        \item[(b)] We say $\calG$ is a \textbf{detecting family} for $\calC$ if the functors $\calC(G,-)$ collectively 'reflect' isomorpisms, i.e. if, given $A \xrightarrow f B$ such that every $G \xrightarrow h B$ with $G \in \calG$ factors uniquely through $f$, then $f$ is an isomorphism.
    \end{enumerate}
    If $\calG = \{G \}$, we call $G$ a \textbf{separator}, or a \textbf{detector} for $\calC$.
\end{fancydefine}
\begin{fancylem}{}\hspace{0.1cm}
    \begin{enumerate}
        \item[(i)] If $\calC$ has equalizers, then any detecting family is separating.
        \item[(ii)] If $\calC$ is balanced, then any separating family is detecting 
    \end{enumerate}
\end{fancylem}
\proof \textit{(i).} Suppose $\calG$ is detecting and suppose $f,g: A \rightrightarrows B$ satisfy the hypothesis of 2.7(a). Then every $G \to A$ with $G \in \calG$ factors uniquely through the equalizer of $f$ and $g$, so the equalizer is an isomorphism. So $f=g$, since an isomorphism is an epimorphism.\gap
\textit{(ii).} Suppose $\calG$ is separating, and $A \xrightarrow B$ satisfies the hypotheses of 2.7(b). If $fg=fh$ for some $g,h : C \rightrightarrows A$, then any $G \xrightarrow k C$ with $G \in \calG$ must satisfy that $gk=hk$, since both are factorizations of $fgk=fhk$ through $f$, so $g=h$, so $f$ is monic. Similarly, if $
l, m: B \rightrightarrows D$ satisfy $lf=mf$, then any $G \xrightarrow n D$ with $G \in \calG$ satisfies $ln =mn$, since it factors through $f$, so $l=m$ so $f$ is epic. Since $\calC$ is balanced, $f$ is an isomorphism. \endproof
\begin{exs} \hspace{0.1cm}
    \begin{enumerate}
    \item[(a)] In \cat{Gp}, $\bbZ$ is a separator and a detector. Similarly, for \cat{Rng}, $\bbZ[X]$ is a separator and a detector.
    \item[(b)] If $\calC$ is small, the set $\{\calC(A,-) \, | \, A \in \ob \calC \} $ is both separating and detecting in $[\calC,\cat{Set}]$ by Yoneda.
    \item[(c)] In \cat{Top}, 1 is a separator, but \cat{Top} has no detecting set of objects.
    \newlec Given an infinite cardinal $\kappa$, let $X_{\kappa}$ be a disrecte space of cardinality $\kappa$, and $Y_{\kappa}$ the same set with the co-$<\kappa$\footnote{Open sets are those whose complements have cardinality $<\kappa$.} topology. The identity $X_{\kappa} \to Y_{\kappa}$ is (continuous but) not a homeomorphism. Given any set $\calG$ of spaces, if $\kappa>\card \calG$ for all $G \in \calG$, then $\calG$ can't detect that $X_{\kappa} \to Y_{\kappa}$ isn't an isomorphism.
    \item[(d)] Let $\calC$ be the category whose objects are ordinals, and in addition to $1_{\alpha} : \alpha \to \alpha$ for each $\alpha$, $\calC$ has two morphisms $f,g: \alpha \rightrightarrows \beta$ whenever $\alpha < \beta$, with composition defined by $ff=fg=gf=gg=f$. $0$ is a detector for $\calC$: it detects that $f,g : 0 \rightrightarrows \alpha$ aren't isomorphisms since neither factors through the other, and it detects that $f,g : \alpha \rightrightarrows \beta$ $(0<\alpha<\beta)$ aren't isomorphisms since $0 \xrightarrow{g} \beta$ doesn't factor through either. But, given any set $\calG$ of ordinals if $\alpha > \gamma$ for all $\gamma \in \calG$, then $\calG$ can't separate $f,g : \alpha \rightrightarrows \alpha+1$.
    \item[(e)] We saw that $\cat{Gp}$ has a single object $\bbZ$ which is both a separator and a detector, but it has no co-separating or co-detecting set of objects. Given any set $\calG$ of groups, let $H$ be a simple group with $\card H > \card G$ for all $G \in \calG$. Then the only homomorphisms from $H$ to members of $\calG$ are trivial, so $\calG$ can't detect that $H \to 1$ isn't an isomorphism.
    \end{enumerate}
\end{exs}
The functors $\calC(A, -) : \calC \to \cat{Set}$ preserve monomorphisms by definition, but not necessarily epimorphisms, so we define 
\begin{fancydefine}{Projective}
    We say an object $P$ of a locally small set $\calC$ is \textbf{projective} if $\calC(P, - )$ preserves epimorphisms, i.e. given
    \begin{center}
    \begin{tikzcd}
        & P \arrow[d,"f"]\\
        Q \arrow[r,"g"] & R
    \end{tikzcd}

\end{center}
there exists $h : P \to Q$ with $gh=f$.\gap If this holds for all $g$ in some class $\calE$ of epimorphisms, we say $P$ is $\calE$-\textbf{projective}. (Dually, $P$ is \textbf{injective} if it's projective in $\calC^{\text{op}}$).
We'll consider the class of pointwise epimorphisms in $[\calC, \cat{Set}]$, i.e. those $\alpha$ such that $\alpha_{A}$ is surjective for all $A$.
\end{fancydefine}
\begin{corr}{}
    Objects of the form $\calC(A, - )$ are pointwise projective in $[\calC, \cat{Set}]$.
\end{corr}
\proof In the diagram of 2.10, if $P=\calC(A,-)$, then $f$ corresponds to some $\Phi(f) \in \calR A$, but $g_{A}$ is surjective, so there exists $\Phi(h) \in QA$ mapping to $\Phi(f)$. \endproof
\begin{fancyprop}{}
    If $\calC$ is small, then $[\calC, \cat{Set}]$ "has enough pointwise projectives," i.e. for any object $F$, there exists a pointwise epimorphism $P \twoheadrightarrow F$ with $P$ pointwise projective.
\end{fancyprop}
\proof Take $P= \bigsqcup_{(A,x)}(\calC(A,-))$ where the disjoint union is over all pairs $(A,x)$ with $A \in \ob \calC$ and $x \in FA.$ Then $P$ is pointwise projective, since the $\calC(A,-)$ are. \gap But we have a natural transformation $\alpha : P \to F$ whose $(A,x)^{\text{th}}$ term\footnote{i.e. the term corresponding to $(A,x)$} is $\Psi(x) : \calC(A, - ) \to F$. And this is pointwise epic, since any $x \in FA$ is in the range of $\Psi(x)$. \endproof









\section{Chapter 3: Adjunctions}
\begin{fancydefine}{Adjunction}
    Let $\calC$ and $\calD$ be categories. By an \textbf{adjunction} between $\calC$ and $\calD$, we mean a pair of functors $F : \calC \to \calD$, $G : \calD \to \calC$, together with a bijection between morphisms $FA \to B$ in $\calD$ and $A \to GB$ in $\calC$ which is natural in both $A$ and $B$. \gap (If $\calC$ and $\calD$ are locally small, this means that $\calD(F-, -)$ and $\calC(-, G-)$ are naturally isomorphic functor $\calC^{\text{op}} \times \calD \to \cat{Set}$.)\gap We say that $F$ is \textbf{left-adjoint} to $G$ , or $G$ is \textbf{right-adjoint} to $F$, and write $(F \dashv G)$.
\end{fancydefine}
\begin{exs}\hspace{0.1cm}
    \begin{enumerate}
        \item[(a)] The free group functor $F : \cat{Set} \to \cat{Gp}$ is left-adjoint to the forgetful functor $U: \cat{Gp} \to \cat{Set}$. Since functions $A \to UG$ correspond bijectively to homomorphisms $FA \to G$, one way comes from this being precisely how we defined the free group functor, the other way is just easy\footnote{Says P.T. Johnstone}.
        \item[(b)] The forgetful functor $\cat{Top} \xrightarrow{U} \cat{Set}$ has a left adjoint $D$ which equips a set with the discrete topology, and has a right-adjoint $I$ which equips a set $A$ with the \textit{indiscrete} topology: $\{\phi, A\}$.
        \item[(c)] The functor $\ob : \cat{Cat} \to \cat{Set}$ has a left-adjoint $D$ sending $A$ to the discrete category with objects $A$. It also has a right-adjoint sending $A$ to the category with objects $A$ and one morphism $a \to b$ for all $a, b \in A$. $D$ also has a left-adjoint, $\pi_{0}$, where $\pi_{0}\calC$ is the quotient of $\ob \calC$ by the smallest equivalence relation identifying $\dom f$ with $\cod f$ for all $f \in \mor \calC$\footnote{Because a discrete category only has objects and identity morphisms.}
        \item[(d)] For any set $A$, we have a functor $(-) \times A : \cat{Set} \to \cat{Set}$ with a right-adjoint given by $\cat{Set}(A, - )$, since functions $B \times A \xrightarrow{f} C$ correspond under under $\lambda$-conversion to functions $B \xrightarrow{\bar f} \cat{Set}(A,C)$. where $\bar f (b)(a) = f(b,a)$. \begin{fancydefine*}{Cartesian Closed}We say a category $\calC$ with binary products defined as in 2.6(f) is \textbf{Cartesian closed} if \[ (-) \times A : \calC \to \calC \] has a right-adjoint $[A,-]$ or $(-)^{A}$ for each $A$.
        \end{fancydefine*}
        E.g. $\cat{Cat}$ is cartesian closed ,with $[\calC,\calD]$ being the usual functor category.
        \item[(e)] \newlec Let $\cat{Idem}$ be the category of pairs $(A,e)$ where $A$ isa set and $e: A \to A$ satisfies $ee=e$. We have a functor $F : \cat{Set} \to \cat{Idem}$ sending $A$ to $(A,1_{A})$ and functions to themselves. We also have a functor $G : \cat{Idem} \to \cat{Set}$ sending $(A,e)$ to the set of fixed points of e: $\{ a \in A \, | \, e(a)=a\}$. And $(F \dashv G)$ since any morphism $FA \to (B,e)$ takes value in $G(B,e)$. But we also have $(G \dashv F)$ since any morphism $(A,e) \xrightarrow{f} FB$ satisfies $f(a)=f(ea)$ for all $a$, so it's determined by its effect on $G(A,e)$. But $G$ isn't faithful, so this isn't an equivalence of categories.
        \item[(f)] Let $\mathbf{1}$ be the discrete category with one object and only one morphism. Then a left adjoint for the unique functor $\calC \to \mathbf{1}$ (if such a functor exists) picks out an \textit{initial object} $I$, i.e. an object for which there's a unique $I \to A$ for any $A \in \ob \calC$. Dually, a right adjoint for $\calC \to \mathbf{1}$ picks out a \textit{terminal object} of $\calC$. Note that, for example, in $\cat{Gp}$, these two coincide, since the trivial group is both initial and terminal. But they don't have to coincide, in $\cat{Set}$, the initial object is the empty set, but the terminal object has one element.
        \item[(g)] Between posets with the inclusion ordering. Given a function $f: A \to B$, if $A' \subseteq A$, and $B' \subseteq B$, then it's easy to see that $Pf(A') \subseteq B'$ if and only if $A' \subseteq P^{*}f(B')$. So $(Pf \dashv P^{*}f)$ as functors between $PA$ and $PB$.
        \item[(h)] Given sets $A$ and $B$ and a relation $R \subseteq A \times B$, we define mappings $ (-)^{r} : PA \to PB$ by $(S)^{r} = \{ b \in B \, | \, (\forall a)(a \in S \Rightarrow (a,b) \in R) \}$ and $(-)^{l} : PB \to PA$ by $(T)^{l}=\{ a \in A \, | \, (\forall b)(b \in T \Rightarrow (a,b) \in R)\}$. These are contravariant functors (= order-reversing mappings), and we have that $S \subseteq T^{l} \Leftrightarrow S \times T \subseteq R \Leftrightarrow T \subseteq S^{r}$. We say $(-)^{l}$ and $(-)^{r}$ are $\textbf{adjoint on the right}$\footnote{This is called a Galois connection, since when we take $A$ to be a field and $B$ to be the group of automorphisms of the field, and the relation $R$ is elements of the field being fixed by elements of the automorphism group, and this gives us a bijection between subfields and subgroups.}
        \item[(i)] The contravariant power-set functor $P^{*}$ is self-adjoint on the right, since functions $A \to PB$ and $B \to PA$ both correspond bijectively to subsets $R \subseteq A \times B$ (i.e. relations).
        \item[(j)] Similarly $(-)^{*} : \cat{Vect}_{k} \to \cat{Vect}_{k}^{\text{op}}$ is self-adjoint on the right, since linear maps $V \to W^{*}$ and $W \to V^{*}$ correspond to bilinear forms on the product $V \times W$. (Need to check naturality with linear maps $V \to V'$ and $W \to W'$ but not hard to do).
    \end{enumerate}
\end{exs}
\begin{fancythm}{}
    Given a functor \[ G : \calD \to \calC\] and $A \in \ob \calC$, let $(A \downarrow G)$ be the category whose objects are pairs $(B,f)$ where $B$ is an object of $\calD$, and $f$ is a morphism $f : A \to GB$ in $\calC$, and whose morphisms $(B,f) \to (B',f')$ are morphisms $B \xrightarrow{g} B'$ in $\calD$ which make this diagram:
    \begin{center}
    \begin{tikzcd}
        A \arrow[r,"f"] \arrow[dr,"f'"] & GB \arrow[d,"Gg"] \\
                                        & GB'
    \end{tikzcd}
    \end{center}
    commute. Then specifying a left-adjoint for $G$ is equivalent to specifying an initial object of $(A \downarrow G)$ for each $A$. (Note: This type of category, $(A \downarrow G)$ is called a comma-category because of the initial (bad) notation $(A,G)$. Now we use a down arrow instead of a comma, since commas can mean anything but the name stuck.)
\end{fancythm}
\proof First suppose we have a left-adjoint $(F \dashv G)$. Then we want to produce an object equipped with a mapping from $A$ to $G$(that object). So we look at $\eta_{A}: A \to GFA$ correspond to $1_{FA}$ under the adjunction, then $(FA,\eta_{A})$ is initial in $(A \downarrow G)$, since if we're given $A \xrightarrow{f} GB$ then \begin{center}
    \begin{tikzcd}
        A \arrow[r, "\eta_{A}"] \arrow[dr, "f"] & GFA \arrow[d, "Gg"]\\ &GB
    \end{tikzcd}
\end{center}
commutes if and only if $g$ is the morphism corresponding to $f$ under the adjunction. So there's a unique such $g$ for any $f$. Therefore $(FA,\eta_{A})$ must be the initial object. 
\gap Conversely, suppose we have initial objects for the categories $(A \downarrow G)$ given by $(FA,\eta_{A})$ for each $A \in \ob \calC$. This defines what $F$ will be on objects (picking out the underlying object of this initial object i.e. $FA$). We make $F$ into a functor by saying that, given $f: A \to A'$, $Ff$ is the unique morphism making \begin{center}
    \begin{tikzcd}
        A \arrow[r, "\eta_{A}"] \arrow[d,"f"] & GFA \arrow[d,"GFf"] \\
        A' \arrow[r,"\eta_{A'}"] & GFA'
    \end{tikzcd}
\end{center}
commute. Functoriality of $F$ follows from uniqueness, if we have $A''$  and a morphism $f'$ $A' \to A''$, then we have $(GFf')(GFf)$ is made equal to $GF(f'f)$ where $f'f$ is the composed morphism from $A$ to $A''$. $F$ is left-adjoint to $G$ since the bijection $A \xrightarrow{f} GB$ and $FA \xrightarrow{g} B$ sends $f$ to the unique $g$ satisfying $(Gg)\eta_{A} = f$. Naturality of the bijection in $A$ was built into the definition of $F$ as a functor, and naturality in $B$ is easy. \endproof This gives us a necessary and sufficient condition for the existence of a left adjoint, which we will be using a lot.
\begin{corr}{}
    If $F$ and $F'$ are both left-adjoint to $G: \calD \to \calC$, then $F \cong F'$ in $[\calC, \calD]$.
\end{corr}
\proof $(FA,\eta_{A})$ and $(F'A, \eta_{A}')$ are both initial in $(A \downarrow G)$. So there's a unique isomorphism $alpha_{A} : (FA, \eta_{A}) \to (F'A, \eta_{A}')$ in this category. And $A \mapsto \alpha_{A}$ is natural, since any morphism from $A$ to $A'$, $f: A \to A'$, we have $alpha_{A'}(Ff)$ and $(F'f)\alpha_{A}$ are both morphisms $(FA, \eta_{A}) \rightrightarrows (F'A', \eta_{A'}'f)$ in $(A \downarrow G)$, so they're equal.
\endproof
\begin{fancylem}{}
    Suppose given \[\calC \mathrel{\mathop{\rightleftarrows}^{\mathrm{F}}_{\mathrm{G}}} \calD \mathrel{\mathop{\rightleftarrows}^{\mathrm{H}}_{\mathrm{K}}} \calE \]
    with $(F \dashv G)$ and $(H \dashv K)$. Then $(HF \dashv GK)$.
\end{fancylem}
\proof We have bijections between morphisms $HFA \to C$. morphisms $FA \to KC$, and then morphisms $A \to GKC$, and they're both natural in $A$ and $C$, and therefore so is their composite. \endproof
\begin{corr}{}
    Suppose given a commuting square 
    \begin{center}
        \begin{tikzcd}
            \calC \arrow[r, "F"] \arrow[d, "G"] & \calD \arrow[d, "H"]\\ \calE \arrow[r, "K"] & \calF
        \end{tikzcd}
    \end{center}
    of categories and functors, such that each of $F,G,H,K$ has a left adjoint. Then the square of left adjoints commutes up to natural isomorphism.
\end{corr}
\proof The two ways round the square are both left-adjoints to $HF=KG$ by (3.5), so (3.4) tells us that they must be naturally isomorphic. \endproof
\newlec Given an adjunction $(F \dashv G)$, we saw in the proof of 3.3 that the morphisms $\eta_{A} : A \to GFA$ corresponding to $1_{FA}$ form a natural transformation $1_{\calC} \to GF$. We call $\eta$ the \textbf{unit} of the adjunction. Dually, we have the \textbf{counit} $\varepsilon: FG \to 1_{\calD}$, where $\varepsilon_{\calB} : FGB \to B$ correspond to $1_{GB}$.
\begin{fancythm}{}
    Given functors \[ F: \calC \to \calD\] and \[G : \calD \to \calC,\] specifying an adjunction $(F \dashv G)$ is equivalent to specifying natural transformations \[ \eta : 1_{\calC} \to GF\] and \[ \varepsilon: FG \to 1_{\calD}.\] satisfying the triangular identity \begin{center}
        \begin{tikzcd}
            F \arrow[r,"F \eta"] \arrow[rd, "1_{F}"]& FGF \arrow[d,"\varepsilon_{F}"]\\
            & F
        \end{tikzcd}
    \end{center}
    \begin{center}
        \begin{tikzcd}
            G \arrow[r, "\eta_{G}"] \arrow[rd, "1_{G}"] & GFG \arrow[d,"G \varepsilon"]\\
            & G
        \end{tikzcd}
    \end{center}
    This is sometimes taken as the \textit{definition} of an adjunction.
\end{fancythm}
\proof \hspace{0.1cm}
\begin{enumerate}
    \item[$\Rightarrow$] Given the adjunction, we've seen how to define $\eta$ and $\varepsilon$, so we need to verify the triangular identity. But $\varepsilon_{FA}$ corresponds to $1_{GFA}$, so by naturality $\varepsilon_{FA}(F\eta_{A})$ corresponds to $1_{GFA} \eta_{A} = \eta_{A}$. But $1_{FA}$ corresponds to $\eta_{A}$, so $\varepsilon_{FA}(F \eta_{A}) = 1_{FA}$. The other identity is dual.
    \item[$\Leftarrow$] Given $\eta$ and $\varepsilon$, we map $A \xrightarrow{f} GB$ to the composite \[ \Phi(f)=FA \xrightarrow{Ff} FGB \xrightarrow{\varepsilon_{B}} B,\] 
    \item[] and $FA \xrightarrow{g} B$ to the composite \[ \Psi(g) = A \xrightarrow{\eta_{A}} GFA \xrightarrow{Gg} GB.\] $\Phi$ and $\Psi$ are natural in $A$ and $B$ since $\eta$ and $\varepsilon$ are natural, so we need to verify that $\Psi \Phi$ and $\Phi \Psi$ are identities. These are dual so we only need verify one. We have $\Psi \Phi(f)$ is the composite: \[ A \xrightarrow{\eta_{A}} GFA \xrightarrow{GFf} GFGB \xrightarrow{G \varepsilon_B} GB\]
    \[ = A \xrightarrow{f} GB \xrightarrow{\eta_GB} GFGB \xrightarrow{G \epsilon_B} GB = A \xrightarrow{f} GB.\]
\end{enumerate}
The other verification is dual. \endproof
In 1.9 we had isomorphisms \[ \alpha : 1_{\calC} \to GF\] and \[ beta : FG \to 1_{\calD}.\] Do these satisfy the triangular identity? Well, no, but they can if we choose correctly.
\begin{fancyprop}{}
    Suppose we had an equivalence $(F,G,\alpha, \beta)$ as in 1.9, then there exist natural isomorphisms \[ \alpha' : 1_{\calC} \to GF\] and \[ \beta' : FG \to 1_{\calD} \] which \textit{do} satisfy the triangular identities. In particular $(F \dashv G)$ and $(G \dashv F)$.
\end{fancyprop}
\proof Actually, we can fix either one of them and only change the other to get the triangular identity. We define $\alpha' = \alpha$, and $\beta'$ to be the composite \[ FG \xrightarrow{(FG \beta)^{-1}} FGFG \xrightarrow{(F \alpha_{G})^{-1}} FG \xrightarrow{\beta} 1_{\calD}.\] Note that $FG \beta = \beta_{FG}$ since \begin{center}
    \begin{tikzcd}
        FGFG \arrow[r, "FG \beta"] \arrow[d, "\beta_{FG}"]& FG \arrow[d, "\beta"]\\
        FG \arrow[r, "\beta"] & 1_{\calD} 
    \end{tikzcd}
\end{center}
commutes by the naturality of $\beta$, and $\beta$ is monic. Dually $GF\alpha = \alpha_{GF}$. For the triangular identities, consider the diagrans
\begin{center}
 \Large COMING SOON \normalsize
\end{center}
where the squares commute by naturality of $\beta^{-1}$ and $\alpha^{-1}$ respectively. So by (3.7) we have $(F \dashv G)$. But $(\beta')^{-1}$ and $(alpha')^{-1}$ also satisfy the triangular identities, so we have $(G \dashv F)$. \endproof
\begin{fancylem}{}
    Suppose given an adjunction $(F : \calC \to \calD \dashv G : \calD \to \calC)$ with counit $\varepsilon: FG \to 1_{\calD}$. Then \begin{enumerate}
        \item[(i)] $\varepsilon$ is pointwise epic $\Leftrightarrow$ $G$ is faithful.
        \item[(ii)] $\varepsilon$ is an isomorphism $\Leftrightarrow$ $G$ is full and faithful  
    \end{enumerate}
\end{fancylem}
\proof \hspace{0.1cm} \begin{enumerate} \item[(i)] Given $B \xrightarrow{g} B'$, the composite $g \varepsilon_{B}$ corresponds under the adjunction to \[Gg : GB \to GB'.\footnote{Pops out of the proof of 3.7 apparently}\] \gap
So $g \mapsto g \varepsilon_{B}$ is injective on morphisms $B \to B'$ iff $G$ is injective. This is true for all $B,B' \Leftrightarrow \varepsilon$ is pointwise epic $\Leftrightarrow G$ is faithful.  

\item[(ii)]Similarly $G$ is full and faithful $\Leftrightarrow g \mapsto g \varepsilon_{B}$ is bijective on morphisms $B \to B'$ for all $B$ and $B'$. This clearly holds if $\varepsilon_{B}$ is an isomorphism for all $B$. \gap

Conversely, if the condition holds, there's a unique $g: B \to FGB$ such that $\varepsilon_{B}g = 1_{B}.$ And $\varepsilon_{B} g \varepsilon_{B} = \varepsilon_{B}$ so $g \varepsilon_{B}$ and $1_{FGB}$ have the same composite with $\varepsilon_{B}$, so they're equal.
\end{enumerate}
\endproof
\begin{fancydefine}{Reflection}
    We call an adjunction satisfying the conditions of 3.9(ii) a \textbf{reflection}. By a \textbf{reflective subcategory} of $\calC$, we mean a full subcategory $\calD$ for which the inclusion has a left adjoint. \gap 
    Note: If $(F: \calC \to \calD \dashv G: \calD \to \calC)$ is a reflection, then by 1.12, $G$ induces an equivalence between $\calD$ and the full subcategory of $\calC$ on objects in the image of $G$, and this subcategory is reflexive.
\end{fancydefine}
\begin{exs}\hspace{0.1cm}
    \begin{enumerate}
        \item[(a)] \cat{AbGp} is reflective in \cat{Gp}. The left adjoint to the inclusion sends $G$\footnote{A group, not a functor.} to its \textit{abelianisation} $\frac{G}{G'}$, where $G'$ is the subgroup of $G$ generated by all commutators $[x,y]=x y x^-1 y^-1$.\gap Note that any homomorphism $G \to A$ where $A$ is abelian factors uniquely through \[ G \to \frac{G}{G'}\footnote{The idea is that we're looking for the \say{best possible approximation} to $G$ which exists in the subcategory, and we have one here which is given by the abelianisation.}.\]
        \item[(b)] Recall that an abelian group $A$ is \textit{torsion} if all its elements have finite order, and \textit{torsion-free} if the only element of finite order is the identity. For a general $A$ the elements of finite order form a subgroup $A_{t}$, and any homomorphism $B \to A$ where $B$ is a torsion group must factor through the inclusion $A_{t} \hookrightarrow A$. So this is a coreflection, and $A \mapsto \frac{A}{A_{t}}$ is a reflection.
        \item[(c)] \newlec The category \cat{KHaus} of compact Hausdorff spaces is reflective in \cat{Top}, the left-adjoint being the \underline{Stone-\v Cech compactification} functor $\beta$. We'll construct $\beta$ in chapter 4.
        \item[(d)] Recall that a subset $C$ of a space $X$ is called \textit{sequentially closed} if $x_{n} \in C$ for all $n$, and $x_{n} \to x$ $\Rightarrow x \in C$. We say $X$ is a \textit{sequential space} if all sequentially closed subsets are closed. The full subcategory \cat{Seq} of sequential spaces if reflective in \cat{Top}: Given $X$, let $X_{s}$ denote the same set, but retopologised so that every sequentially closed subset is closed.The identity mapping from $X_{s} \to X$ is the counit of the adjunction.
        \item[(e)] The category \cat{Preord} of small preorders is reflective in \cat{Cat}. The reflector sends $\calC$ to $\frac{\calC}{\simeq}$, where $\simeq$ is the congruence identifying all paralle pairs in $\calC$.
        \item[(f)] For any space $X$, the poset $\Omega(X)$ of open subsets of $X$ is corefelctive in $PX$. Since if $U$ is open and $A \subseteq X$ is arbitrary, then $U \subseteq A$ if and only if $U \subseteq A^{0}$ (so $(-)^{0}$ is right-adjoint to $\Omega(X) \hookrightarrow PX$). Dually, the poset of closed subsets is reflective in $PX$. 
    \end{enumerate}
\end{exs}
\section{Limits and Colimits}
To define (co)limits, we need a formal definition of a diagram in a category.
\begin{fancydefine}{Diagram}
    Let $J$ be a category (almost always small, often finite). By a \textbf{diagram} of shape $J$ in $\calC$, we mean a functor \[ D : J \to \calC. \] We call the objects $D(j), j \in \ob J$, \textbf{vertices} of $D$. and the morphisms $D(\alpha)$, $\alpha \in \mor J$, \textbf{edges} of $D$.
\end{fancydefine}
For example, let $J$ be the category 
\begin{center}
    \begin{tikzcd}
        \cdot \arrow[r] \arrow[d] \arrow[rd] & \cdot \arrow[d] \\ \cdot \arrow[r] & \cdot 
    \end{tikzcd}
\end{center}
A diagram of shape $J$ in $\calC$ is a commutative square. If $J'$ is 
\begin{center}
    \begin{tikzcd}
        \cdot \arrow[r] \arrow[d] \arrow[rd, shift right =0.5ex] \arrow[rd, shift left=0.5ex] & \cdot \arrow[d] \\ \cdot \arrow[r] & \cdot 
    \end{tikzcd}
\end{center}
Then a diagram of shape $J'$ is a not-necessarily-commutative square.
\begin{fancydefine}
    Given a diagram $D : J \to \calC$, a \textbf{Cone over} $\mathbf{D}$ consists of an object $A$ (its \textbf{apex}) together with morphisms $\lambda_{j} : A \to D(j)$ its \textbf{legs} for each $j \in \ob J$, such that \begin{center}
        \begin{tikzcd}
            & A \arrow[dl, "\lambda_j"] \arrow[dr, "\lambda_j'"] &\\
            D(j) \arrow[rr, "D(\alpha)"] & & D(j')
        \end{tikzcd}
    \end{center}
    commutes for each $\alpha \in \mor J$. So it basically overall will look like a cone.
    \gap
    \begin{define*}{Morhism of Cones/Limit}
    A \textbf{morphism of cones} $(A, (\lambda_{j} \, | \, j \in \ob J)) \to (B, (\mu_{j} \, | \, j \in \ob J))$ is a morphism $A \xrightarrow{f} B$ such that \begin{center} \begin{tikzcd}
        A \arrow[dr, "\lambda_j"] \arrow[rr, "f"] & & B \arrow[dl,"\mu_j"]\\
        & D(j) &
    \end{tikzcd}
    \end{center} commutes for all $j$. We write $\Cone (D)$ for the category of cones over $D$. A \textbf{Limit} for $D$ is a terminal object of $\Cone (D)$, if this exists. Dually, we have the notion of a \textbf{cone under} a diagram, and a \textbf{colimit} for $D$ is an \textit{initial} cone under $D$.\gap 
    Note that a cone over $D$ with apex $A$ is the same thing as a natural transformation \[ \Delta A \to D \] where $\Delta A$ is the \textbf{Constant diagram} with all vertices $A$ and all edges $1_{A}$. \gap 
    $\Delta$ is a functor $\calC \to [J,\calC]$ and $\Cone (D)$ is the arrow category $(\Delta, \downarrow, D)$ defined as in (3.3). \gap 
    So $\calC$ has limits for all diagrams of shape $J$ $\Leftrightarrow$ $\Delta : \calC \to [J, \calC]$ has a right adjoint, by (3.3).
\end{define*}
\end{fancydefine}
\begin{exs}\hspace{0.1cm}
    \begin{enumerate}
        \item[(a)] If $J = \emptyset$, there's a unique diagram of shape $J$ in any category. This diagram is drawn below: \gap \hspace{0.1cm} \gap \hspace{0.1cm} \gap \hspace{0.1cm} \gap \hspace{0.1cm} \gap A cone over it is just an object of $\calC$ (the apex), so $\Cone(D) \cong \calC$. So a limit for $D$ is a terminal object for $\calC$.
        \item[(b)] If $J$ is the discrete category with two objects, a diagram of shape $J$ is a pair of objects $A,B$. A cone over it is a \textit{span:} \begin{center}
            \begin{tikzcd}[column sep = small]
                & C \arrow[dl] \arrow[dr] &\\
                A & & B
            \end{tikzcd}
        \end{center}
        A limit cone is a categorical product:
        \begin{center}
            \begin{tikzcd}[column sep = small]
                & A \times B \arrow[dl, "\pi_1"] \arrow[dr, "\pi_2"] &\\
                A & & B
            \end{tikzcd}
        \end{center}
        as defined in 2.6(f). Dually a colimit for $A,B$ is a coproduct $A+B$.
        \item[(c)] IF $J$ is any discrete category, a diagram of shape $J$ is a family $(A_{j}| j \in J)$ of objects, and a (co)limit for it is a (co)product: \begin{align*}
            \Pi_{j \in J} A_{j} & & \sum_{j \in J} A_{j}
        \end{align*}
        \item[(d)] If $J$ is \[\cdot \rightrightarrows \cdot,\] a diagram of shape $J$ is a parallel pair $f,g : A \rightrightarrows B$. A cone over this diagram consists of 
        \begin{center}
            \begin{tikzcd}[column sep = small]
                & C \arrow[dl, "h"] \arrow[dr, "k"] &\\
                A & & B
            \end{tikzcd}
        \end{center}
        satisfying $fh=k=gh$, or equivalently a morphism $C \xrightarrow{h} A$ satisfying $fh=gh$. So a limit for $f,g : A \rightrightarrows B$ is an equalizer for $(f,g)$ as defined in 2.6(g) (and a colimit for it is a coequalizer).
        \item[(e)] Let $J$ be the category \begin{center}
            \begin{tikzcd}
                & \cdot \arrow[d]\\
                \cdot \arrow[r] & \cdot
            \end{tikzcd}
        \end{center}
        So a diagram of shape $J$ is a \textit{cospan}:
        \begin{center}
            \begin{tikzcd}
                & A \arrow[d, "f"]\\
                B \arrow[r, "g"] & C
            \end{tikzcd}
        \end{center}
        A cone over it has three legs\begin{center} \begin{tikzcd}
            D \arrow[r,"h"] \arrow[d, "k"] \arrow[dr, "e"] & A\\
            B & C
        \end{tikzcd}
    \end{center}
        but $e$ is redundant, since $e=fh=gk$. So it's a span completing $(f,g)$ to a commutative square. A limit is called a \textit{pullback} of $(f,g)$. Dually, colimits for diagrams of shape \begin{center}
            \begin{tikzcd}
                \cdot \arrow[r] \arrow[d] & \cdot\\
                \cdot
            \end{tikzcd}
        \end{center}
            are called \textit{pushouts}. \newlec In any category with binary products and equalizers, we can construct a pullback for \begin{center}
                \begin{tikzcd}
                    & A \arrow[d,"f"]\\
                    B \arrow[r,"g"] & C\\
                \end{tikzcd}
            \end{center}
            by first forming the product $A \times B$ and then forming the equalizer of $f\pi_{1},g\pi_{2} : A \times B \rightrightarrows C$.
        \item[(f)] Let $M$ be the 2-element monoid $\{1,e\}$ with $e^{2}=e$. A diagram of shape $M$ is an object equipped with an idempotent endomorphism. A cone over it is $B \xrightarrow{f} A$ satisfying $ef=f$. So a limit (resp. colimit) for this diagram is the monic (resp. epic) part of a splitting of $e$ by $2.14$. (I.e. the maps to $A$ to the object forming the cone over this diagram split $e$). Note that in 3.2(e) the functor $F : \cat{Set} \to [M, \cat{Set}]$ is just $\Delta$. So this explains why its left and right adjoints coincide.
        \item[(g)] Let $\bbN$ be the ordered set of natural numbers. A diagram of shape $\bbN$ is a \textbf{direct sequence} \[ A_{0} \to A_{1} \to A_{2} \to A_{3} \to A_{4} \to A_{5} \to A_{6} \to A_{7} \to A_{8} \to A_{9} \to A_{10} \to A_{11} \to A_{12} \to A_{13} \to A_{14} \to A_{15} \to A_{16} \to A_{17} \to A_{18} \to \]
        A colimit for it is a \textbf{direct limit} $A_{\infty}$ equipped with morphisms $A_{n} \to A_{\infty}$ for all $n$ compatible with the transition morphisms $A_{n} \to A_{n+1}$ and universal among such. \gap Dually, an \textbf{inverse sequence} is a diagram of shape $\bbN^{\text{op}}$, and a limit for it is an \textbf{inverse limit}. For example, an infinite dimensional CW-complex $X$ can be viewed as the direct limit of its $n$-dimensional skeletons $X_{n}$. \gap Also, the ring of $p$-adic integers is the limit of the inverse sequence of rings \[ \dots \to \frac{\bbZ}{p^{2}\bbZ} \to \frac{\bbZ}{p\bbZ} \to \frac{\bbZ}{\bbZ}. \]
    \end{enumerate}
\end{exs}
\begin{fancylem}{}
    \begin{enumerate}
        \item[(i)] If $\calC$ has equalizers and small products, then $\calC$ has all small limits.
        \item[(ii)] If $\calC$ has equalizers and all finite products, then $\calC$ has all finite limits.
        \item[(iii)] If $\calC$ has pullbacks and a terminal object, then $\calC$ has all finite limits.
    \end{enumerate}
\end{fancylem}
\proof \hspace{0.1cm}\begin{enumerate}
    \item[(i) \& (ii)] Suppose given $D: J \to \calC$. Form the products $P = \Pi_{j \in \ob J} D(j)$ and $Q = \Pi_{\alpha \in \mor \calC} D(\cod \alpha)$. We have morphisms \[ f,g : P \rightrightarrows Q\] defined by $\pi_{\alpha}f = \pi_{\cod \alpha}$ and $\pi_{\alpha}g=D(\alpha) \pi_{\dom \alpha}.$ Let $E \xrightarrow{e} P$ be an equalizer for $(f,g)$ and define $\lambda_{j} = \pi_{J}e : E \to D(j)$. 
    \gap These both form a cone over $D$, since if  $\alpha : j \to j'$, we have \[D(\alpha)\lambda_{j} = D(\alpha)\pi_{j}e = \pi_{\alpha}ge = \pi_{\alpha}fe = \pi_{j'}e = \lambda_{j'}.\] Then given any cone $(A, (\mu_{j} | j \in \ob J))$ over $D$, we have a unique $\mu : A \to P$ with $\pi_{j} \mu = \mu_{j}$ with $\pi_{j}\mu = \mu_{j}$ for all $j$. Then \[ \pi_{\alpha}f\mu = \mu_{\cod \alpha} = D(\alpha) \mu_{\dom \alpha} = \pi_{\alpha} g \mu\] for all $\alpha$, so $\mu$ factors uniquely through $e$.
    \item[(iii) $\Rightarrow$ (ii)] If $1$ is terminal in $\calC$, the pullback of \begin{center}
        \begin{tikzcd}
            & A \arrow[d]\\
            B \arrow[r] & 1\\
        \end{tikzcd}
    \end{center}
    has the universal property of a product on $A\times B$. Then we can form the product of any finite number of objects $\Pi_{i = 1}^{n} A_{j}$ inductively. Then all that is left is to form the equalizer for $f, g : A \rightrightarrows B$. Consider the pullback of \begin{center}
        \begin{tikzcd}
            & A \arrow[d, "(1_{A} {,} f)"]\\
            A \arrow[r, "(1_{A} {,} g)"] & A \times B
        \end{tikzcd}
    \end{center}
    Any conve over this has its two legs $C \rightrightarrows A$ being equal, so a pullbacj for it is also an equalizer for $(f,g)$.
\end{enumerate}
\endproof
\begin{fancydefine*}{Complete}
    We say $\calC$ is \textbf{(co)complete} if it has all small (co)limits.
\end{fancydefine*}
We know that $\cat{Set}$, $\cat{Gp}$, \cat{Top} are complete and cocomplete.
\begin{fancydefine}{}\hspace{0.1cm}
    \begin{enumerate}
    \item[(a)] We say $G$ \textbf{preserves} limitsof shape $J$ if, given $D: J \to \calD$, and a limit cone $(L, (\lambda_{j} \, | \, j \in \ob J))$ for $D$, the cone $(GL, (G\lambda_{j} \, | \, j \in \ob J))$ is a limit for $GD$.
    \item[(b)] We say $G$ \textbf{reflects} limits of shape $J$ if, given $D : J \to \calD$, and a cone $(L, (\lambda j))$ such that $(GL, (G\lambda_{j}))$ is a limit for $GD$, then $(L, (\lambda_{j}))$ is a limit for $D$.
    \item[(c)] We say $G$ \textbf{creates} limits of shape $J$ if, given $D : J \to \calD$ and a limit $(M, (\mu_{j}))$ for $GD$ in $\calC$, there exists a cone $(L, (\lambda_{j}))$ over $D$ such that $(GL, (G\lambda_{j})) \cong (M, (\mu_{j}))$ in $\Cone(GD)$ and any such conve is a limit for $D$\footnote{There is another common definition which requires a unique lifting to limits.}.
    \end{enumerate}
    Normally, in $(a)$ we assume that $\calD$ has all limits of shape $J$, and in $(b)$ and $(c)$ we assume $\calC$ has them. With these assumptions \[ \text{$G$ creates limits of shape $J$} \Leftrightarrow \text{$G$ preserves and reflects them, and $\calD$ has them.} \]
\end{fancydefine}
\begin{corr}{}
    In any of the statements of 4.4, we can replace both instances of $\calC$ has by either \say{$\calD$ has and $G: \calD \to \calC$ preserves} or \say{$\calC$ has and $G: \calD \to \calC$ creates.}
\end{corr}
\proof Omitted. Not hard.
\begin{exs}\hspace{0.1cm}
    \begin{enumerate}
        \item[(a)] The forgetful functor $U$ : \cat{Gp} $\to$ \cat{Set} creates all small limits. It doesn't preserve coproducts or (therefore) colimits.
        \item[(b)] The forgetful functor $U : \cat{Top} \to \cat{Set}$ preserves all small limits and colimits.
        \item[(c)] The inclusion $\cat{AbGp} \to \cat{Gp}$ reflects coproducts, but doesn't preserve them (since products in \cat{AbGp} are direct sums, but in \cat{Gp} they are free groups (which are \textit{never} abelian (unless one summand is trivial))). But the coproduct of ${1}$ and $H$ is $H$ in both categories.
        \item[(d)] Suppose \newlec $\calD$ has limits of shape $J$. Then, for any $\calC$, the functor category $[\calC, \calD]$ has limits of shape $J$, and the forgetful functor $[\calC, \calD] \to \calD^{\ob \calC}$ creates them. Given a diagram $D: J \to [\calC, \calD]$, we can regard it as a functor $J \times \calC \xrightarrow{\calD} \calD$; for each $A \in \ob \calC$, $\calD( - , A)$ is a diagram of shape $J$ in $\calD$, and so has a limit $(LA, (\lambda_{j,A} | j \in \ob J))$. Given any $f: A \to B$ in $\calC$, the composites $LA \xrightarrow{\lambda_{j,A}} \calD(j,A) \xrightarrow{\calD(j,f)} \calD(j,B)$ form a cone over $\calD(-, B)$, and so induce a unique $Lf : LA \to LB$. As usual, uniqueness implies that $L$ is a functor $\calC \to \calD$, and it's the unique lifting of $(LA, A \in \ob \calC)$ to an object of $[\calC, \calD]$ which makes the $\lambda_{j,A}$ into natural transformations. This is a limit conve in $[\calC, \calD]$ because given any cone $(M, (\mu_{j, - } | j \in \ob J))$ over $\calD$ in $[\calC, \calD]$, the $\mu_{j,A}$ form a cone over $\calD(-,A)$, so we get a unique $\nu_{A} : MA \to LA$ such that $\lambda_{j,A}\nu_{A} = \mu_{j,A}$ for all $A$. And the $\nu$'s form a natural transformation $M \to L$ because given any $f: A \to B$ in $\calC$, $\nu_{B}(Mf)$ and $(Lf)\nu_{A}$ are factorizations of the same cone over $\calD(-, B)$ through its limits.
    \end{enumerate}
\end{exs}
\begin{remark}
    Note that $A \xrightarrow{f} B$ is monic if and only if \begin{center}
        \begin{tikzcd}
            A \arrow[r, "1_A"] \arrow[d, "1_A"] & A \arrow[d, "f"]\\
            A \arrow[r, "f"] & B
        \end{tikzcd}
    \end{center}
    is a pullback. Hence if $\calD$ has pullbacks, any monomorphism in $[\calC, \calD]$ is a pointwise monomorphism, since this pullback must be constructed pointwise by 4.7(d).
\end{remark}
\begin{fancylem}{}
    Suppose $G : \calD \to \calC$ has a left adjoint. Then $G$ preserves all limits which exist in $\calD$.
\end{fancylem}
\proof We give two proofs for this, the second proves the claim and the first indicates why we should believe the claim. \gap \textit{Proof 1.} Suppose $\calC$ and $\calD$ both have limits of shape $J$. If $(F \dashv G)$, then the diagram:
\begin{center}
    \begin{tikzcd}
        \calC \arrow[r, "F"] \arrow[d, "\Delta"] & \calD \arrow[d, "\Delta"]\\
        {[J,\calC]} \arrow[r, "{[J , F]}"] & {[J,\calD]}
    \end{tikzcd}
\end{center}
commutes, But all functors in this diagram have right adjoints, so by (3.6), we know that 
\begin{center}
    \begin{tikzcd}
        {[J,\calD]} \arrow[r, "{[J,G]}"] \arrow[d, "\lim_{J}"] & {[J, \calC]} \arrow[d,"\lim_{J}"]\\
        \calD \arrow[r,"G"] & \calC
    \end{tikzcd}
\end{center}
commutes up to isomorphism. So $G$ preserves limits of shape $J$. \endproof 
\textit{Proof 2.} Suppose we are given $D: J \to \calD$ and a limit con $(L, (\lambda_{j} \, | \, j \in \ob J))$ over it. Given a cone over $GD$ with apex $A$ and legs $\mu_{j} : A \to GD(j)$, the $\mu_{j}$ correspond to morphisms $\bar \mu_{j}: FA \to D(j)$, which form a cone over $D$ by naturality of the adjunction. So we get a unique $\bar \mu : FA \to L$ with $\lambda_{j} \bar mu = \bar \mu_{j}$ for all $j$. That is to say, such that $(G \lambda_{j})\mu = \mu_{j}$ for all $j$, where $\mu : A \to GL$ corresponds to $\bar \mu$. \endproof
Supposing that $\calD$ has all limits, and $G :\ calD \to \calC$ preserves all limits, we'd expect $G$ to have a left adjoint.
\begin{fancylem}{}
    Suppose $\calD$ has all limits and $G : \calD \to \calC$ preserves limits of shape $J$. Then for any $A \in \ob \calC$, $(A \downarrow G)$ has limits of shape $J$, and the forgetful functor $U : (A \downarrow G) \to \calD$ creates them.
\end{fancylem}
\proof Suppose we are given $D : J \to (A \downarrow G)$, write each $D(j)$ as $(UD(j), f_{j})$ where $f_{j} : A \to G \cup D(j)$. Let $(L,(\lambda_{j} \, | \, j \in \ob J))$ be a limit for $UD$ in $\calD$, then $(GL, (G\lambda_{j} \, | \, j \in \ob J))$ is a limit for $GUD$. But the edges of $D$ are morphisms in $(A \downarrow G)$, so the $f_{j}$ form a cone over $GUD$. So we get a unique $f: A \to GL$ such that $(G \lambda_{j})f = f-{j}$ for all $j$. In other words, we have a unique lifting of $L$ to an object $(L,f)$ of $(A \downarrow G)$ which makes the $\lambda_{j}$'s into a cone over $D$ with apex $(L,f)$. Given any cone over $D$ with apex $(M,g)$, it's a cone over $UD$ with apex $M$, and so we get a unique $h: M \to L$; this is a morphism in $(A \downarrow G)$ since both $(Gh)g$ and $f$ are factorizations through $L$ of the same cone over $UD$. \endproof
\begin{fancylem}{}
    Let $\calC$ be a category. Specifying an initial object of $\calC$ is equivalent to specifying a limit for $1_{\calC} : \calC \to \calC$.
\end{fancylem}
\proof First suppose $I$ is an initial object. Then the unique morphisms $I \to A$ form a cone over $1_{\calC}$. If $(B, (\lambda_{A} : B \to A))$ is any other cone, then $\lambda_{I}$ is its unique factorization through this one. \gap Conversely, suppose $(I, (\lambda_{A} : I \to A \, | \, A \in \ob \calC))$ is a limit for $1_{\calC}$. Then $I$ is weakly initial. If $f : I \to A$ is any morphism, then $f \lambda_{I}=\lambda_{A}$, so we need to show that $\lambda_{I}$ is the identity morphism on $I$. But in particular, $\lambda_{A} \lambda_{I}=\lambda_{A}$ for all $A$, so $\lambda_{I}$ is a factorisation of the limit cone through itself. But the limit cone has to map uniquely through itself, so $\lambda_{I}$ must be $1_{I}$. \endproof
These last two lemmas taken together prove the \say{Primitive Adjoint Functor Theorem:} If $\calD$ has all limits and $G : \calD \to \calC$ preserves all limits, then each $(A \downarrow G)$ has all limits by 4.10, so they have initial objects by 4.11, so $G$ has a left-adjoint.
\begin{fancythm}{General AFT} \newlec Suppose $\calD$ is complete and locally small. Then $G: \calD \to \calC$ has a left-adjoint if and only if $G$ preserves small limits and satisfies the \say{Solution-set condition}, which is that:
    For any object $A \in \ob \calC$, there's a set $\{ f_{i} : A \to GB_{i} \, | \, i \in I\}$, such that every $f: A \to GB$ factors as \[ A \xrightarrow{f_{i}} GB_{i} \xrightarrow{Gg} GB\] for some $i$ and some $g : B_{i} \to B$.
\end{fancythm}
\proof \hspace{0.1cm}
\begin{enumerate}
    \item[$\Rightarrow$] If $(F \dashv G)$, $G$ preserves limits by (4.9), by (3.3) we know that the category $(A \downarrow G)$ has an initial object $\{\eta_{A} : A \to GFA\}$, which is a solution set at $A$, for any $A$.
    \item[$\Leftarrow$] Fix $A \in \ob \calC$, by 4.10, $(A \downarrow G)$ is complete, and is locally small since morphisms $(B,f) \to (B',f')$ in $(A \downarrow G)$ are a subset of $\calD(B,B')$. So we need to show that if $\scrA := (A \downarrow G)$ is complete and locally small, and has a weakly initial set of objects $\{S_{i} \, | \, i \in I\}$, then it has an initial object. \gap First, we form $P = \Pi_{i \in I}S_{i}$, which must be a weakly initial object, since we have morphisms $P \xrightarrow{\pi_{i}} S_{i}$ for all $i$. Now consider the diagram \begin{center}
        \begin{tikzcd}
        P \arrow[r, yshift=1ex] \arrow[r, yshift=0.5ex] \arrow[r] \arrow[r, yshift=-0.5ex] \arrow[r, yshift=-1ex] & P
        \end{tikzcd}
    \end{center}
    whose edges are all endomorphisms of $P$ (which is still a set-sized diagram, by local smallness). This diagram must have a limit by assumption, so let $I \xrightarrow{i} P$ be the limit for this diagram. Then $I$ is weakly initial. To show $I$ is initial, suppose there were two morphisms 
    \begin{tikzcd}
        I \arrow[r, yshift=0.5ex, "f"] \arrow[r, yshift=-0.5ex, "g", swap] & C
    \end{tikzcd}, form their equalizer $E \xrightarrow{e} I$ and choose some $P \xrightarrow{h} E$. Now $ieh$ and $1_{P}$ are endomorphisms of $P$, so $iehi = 1_{P}i=i$. But $i$ is monic, so we get $ehi=1_{I}$. But then $e$ must be a split epimorphism, and hence $f=g$.
\end{enumerate}
\endproof
\begin{exs}{}
    \item[(a)] Consider the forgetful functor $U: \cat{Gp} \to \cat{Set}$, $\cat{Gp}$ is complete and preserves small limits (in fact it creates them, by 4.7(a)), and $\cat{Gp}$ is locally small. Given a set $A$, any function $f: A \to UG$ can be factors as \[ A \rightarrow UG' \rightarrow UG,\] where $G'$ is the subgroup generated by $\{f(a) \, | \, a \in A\}$. And $\card(G') \leq \max(\aleph_{0}, \card A)$, sow e fix a set $B$ of this cardinality. Consider all group structures on subsets $B' \subseteq B$, and all possible functions $A \to B'$, and these form a solution-set at $A$. \gap \textbf{Note.} Cardinality bound on $G'$ contains (most of) what we need to construct free groups explicitly.
    \item[(b)] Let \cat{CLat} be the category of complete lattices (=posets with all meets and joins). The forgetful functor $\cat{CLat} \to \cat{Set}$ creates limits just like $\cat{Gp} \to \cat{Set}$. But Andrew Hales (1969) proved that there are arbitrarily large 3-generator complete lattices. So $U$ has no solution-set at $A=\{a,b,c\}$. However, note that $U$ \textit{is} representable, so $(1 \downarrow U)$ has an initial object. But \cat{CLat} doesn't have coproducts (not even finite ones), so this doesn't yield initial objects for $(A \downarrow U)$.
\end{exs}
\begin{fancydefine}{subobject}
    By a \textbf{subobject} of an object $A$, we mean a monomorphism \begin{center}\begin{tikzcd}
        A \arrow[r, rightarrowtail, "m"] & A
    \end{tikzcd}
\end{center}
The subobjects of $A$ in $\calC$ form a preorder $\Sub_{\calC}(A)$ with $m \leq m'$ if and only if $m$ factors through $m'$. We say $\calC$ is \textbf{well-powered} if $\Sub_{\calC}(A)$ is equivalent to a (small) poset for any $A$. Dually $\calC$ can be \textbf{well-copowered} (not cowell-powered\footnote{Not powered by Simon Cowell.})
\end{fancydefine}
For example $\cat{Set}$ is well-powered, since every $A' \rightarrowtail A$ is isomorphic to a subset inclusion; \cat{Set} is also well-copowered, since quotients correspond (up to isomorphism) to equivalence relations on $A$.
\begin{fancylem}{}
    Given a pullback square \begin{center}
        \begin{tikzcd}
                P \arrow[d, "k"] \arrow[r, "h"] & A \arrow[d,rightarrowtail, "f"]\\
                B \arrow[r, "g"] & C
        \end{tikzcd}
    \end{center}
    with $f$ monic, then $k$ is also monic. (Monomorphisms are stable under pullbacks).
\end{fancylem}
\proof Suppose we are given \begin{tikzcd} D \arrow[r, yshift=0.5ex, "l"] \arrow[r, yshift=-0.5ex, "m", swap] & P \end{tikzcd} with $kl=km$,. Then $fhl=gkl =gkm =fhm$, but $f$ is monomorphism, so $hl=hm$. So $l$ and $m$ are both factorisations of $(hl,kl)$ through the pullback, so $l=m$. \endproof
\begin{fancythm}{Special Adjoint Functor Theorem}
    Suppose $\calC$ and $\calD$ are both locally small, and that $\calD$ is complete and well-powered, and $\calD$ has a coseparating set of objects. (See 2.7(a)). Then a functor $G: \calD \to \calC$ has a left-adjoint if and only if $G$ preserves small limits.
\end{fancythm}
\proof\hspace{0.1cm} \begin{enumerate}
    \item[$\Rightarrow$] This is easy, this just comes from 4.9.
    \item[$\Leftarrow$] Fix $A \in \ob \calC$. As before $(A \downarrow G)$ inherits completeness and local smallness from $\calD$. It's also well-powered, since subobjects of $(B,f)$ in $(A \downarrow G)$ are just those \begin{tikzcd} B' \arrow[r, rightarrowtail, "m"] & B \end{tikzcd} in $\calD$ for which $f$ factors through $Gm$. (Note that $(A \downarrow G)$ preserves monomorphisms, since it preserves pullbacks.) \gap 
    And if $\{S_{i} \, | \, i \in I\}$ is a coseparating set for $\calD$, then the set of all $f: A \to GS_{i}$ with $i \in I$ is coseparating in $(A \downarrow G)$, since given \begin{tikzcd} {(B,g)} \arrow[r,yshift=0.5ex,"h"] \arrow[r,yshift=-0.5ex,"k", swap] & {(B',g')} \end{tikzcd} with $h \neq k$, we can find $B' \xrightarrow{l} S_{i}$ with $lh \neq lk$ and $l$ is a morphism $(B',g') \to (S_{i} \, | \, (Gl)g')$ in $(A \downarrow G)$. \gap It \newlec remains to show that if $\calA$ is complete, locally small, and well-powered, and has a coseparating set \[ \{S_{i} \, | \, i \in I \}, \] then it has an initial object. First form $P = \Pi_{i \in I}S_{i}$. Consider the diagram \begin{center}\begin{tikzcd}
        \cdot \arrow[dr, rightarrowtail] & \cdot \arrow[d,rightarrowtail]\\
        \cdot \arrow[r,rightarrowtail] & P \end{tikzcd} \end{center}
        (but with way more objects mapping to $P$) whose edges are the members of a representative set of subobjects of $P$. If $I$ is the apex of a limit cone for this, the legs of the limit cone are all monic by the argument of (4.15), so in particular $I \to P$ is monic, so it's a least subobject of $P$. IF we had \begin{tikzcd} I \arrow[r, yshift=0.5ex, "f"] \arrow[r, yshift=-0.5ex, "g", swap] & A \end{tikzcd}, their equalizer $E \rightarrowtail I$ would yield a subobject of $P$ contained in $I \rightarrowtail P$, so $E \rightarrowtail I$ would be an isomorphism. \gap Given any object $A$ of $\calA$, form the product $Q = \prod_{(i,f)} S_{i}$ indexed over pairs $(i,f)$ with $i \in I$ and $f:A \to S_{i}$, and define $g: A \to Q$ by $\pi_{(i,f)}g=f$. Since the $S_{i}$ are coseparating, $g$ is monic. We also have $h :  P \to Q$ defined by $\pi_{(i,f)}h=\pi_{i}$, so in the pullback: \begin{center} \begin{tikzcd}
            B \arrow[r] \arrow[d, rightarrowtail, "k"] & A \arrow[d, rightarrowtail, "g"]\\
            P \arrow[r, "h"] & Q
        \end{tikzcd}
    \end{center}
    $k$ is monic by (4.15), so it's a subobject of $P$, so we have a factorisation \begin{center}
        \begin{tikzcd}
            I \arrow[r, dashrightarrow], \arrow[dr, rightarrowtail] & B \arrow[d, rightarrowtail, "k"]\\
            & P
        \end{tikzcd}
    \end{center}
    and hence a morphism $I \to A$.
\end{enumerate}
\endproof
\begin{ex}
    Consider the inclusion functor $I: \cat{KHaus} \to \cat{Top}.$ \cat{KHaus} is closed under products in \cat{Top} by Tychonoff's theorem, and it's closed under equalizers, since the equalizer of 
    \begin{tikzcd}
         X \arrow[r, "f", yshift=0.5ex] \arrow[r,"g", yshift=-0.5ex, swap] & Y 
        \end{tikzcd} 
        with $Y$ Hausdorff is a closed subspace of $X$, and hence compact if $X$ is. Therefore, by (4.4(i)) \cat{KHaus} is complete, and $I$ preserves small limits.\gap \cat{KHaus} is locally small, and well-powered since subobjects of $X$ are isomorphic to closed subspaces. And \cat{KHaus} has a coseparator $[0,1]$ by Urysohn's Lemma. So by (4.16), $I$ has a left-adjoint.
\end{ex}

\begin{remarks} \hspace{0.1cm}
    \begin{enumerate}
        \item[(a)] Cech's construction of $\beta$ was virtually identical to that provided by (4.17). Given $X$, he formed \[ P = \prod_{f:X \to [0,1]} [0,1]\] and $g: X \to P$, defined by $\pi_{f}g=f$. That is to say, he formed the product of all the members of a coseparating set for $(X \downarrow I)$. He then defined $\beta X$ to be the closure of the image of $g$, i.e. the smallest subobject of $(P,g)$ in $(X \downarrow I)$. So (4.16) may have been inspired by Cech's proof.
        \item[(b)] We could construct $\beta$ using the general adjoint functor theorem. To get a solution-set at a space $X$, observe that any morphism from $X$ to a compact hausdorff space $f: X \to IY$ factors as \[ X \to IY' \to IY\] where $Y'$ is the closure of $X' = \{f(x) \, | \, x \in X \}$. And, if $Y'$ is Hausdorff and $X'$ is dense in it, then $\card Y' \leq 2^{2^{\card X'}}$.
    \end{enumerate}
\end{remarks}
\section{Monads}
Suppose I have an adjunction \begin{tikzcd} \calC \arrow[r, "F",yshift=0.5ex] & \calD \arrow[l, "G", yshift=-0.5ex] \end{tikzcd} where $\calC$ is familiar, but $\calD$ isn't. How much of this adjunction can we \say{see} in terms of $\calC$.\gap Well, we have the composite $T=GF: \calC \to \calC$, and the unit $\eta : 1_{\calC} \to T$, and also the natural transformtion $\mu = G \varepsilon_{F} : TT \to T$. The triangular identities of 3.7 yield the commutative diagrams \begin{enumerate} \item \begin{center}
    \begin{tikzcd}
        T \arrow[r, "T \eta"] \arrow[dr, "1_T"] & TT \arrow[d, "\mu"]\\
        & T
    \end{tikzcd}
\end{center}
\item \begin{center}
    \begin{tikzcd}
        T \arrow[r, "\eta_T"] \arrow[dr, "1_T"] & TT \arrow[d, "\mu"]\\
         & T
    \end{tikzcd}
\end{center}
and naturality of $\varepsilon$ yields \item \begin{center}\begin{tikzcd}
    TTT \arrow[d, "\mu_{T}"] \arrow[r,"T_{\mu}"] & TT \arrow[d,"\mu"]\\
    TT \arrow[r, "\mu"] & T
\end{tikzcd}
\end{center}
\end{enumerate}
\begin{fancydefine}{Monad}
    By a \textbf{monad} on a category $\calC$, we mean a triple $\bbT= (T,\eta, \mu)$ where $T$ is a functor $\calC \to \calC$ and $\eta: 1_{\calC} \to T$, $\mu: TT \to T$ are natural transformations satisfying 1-3 above.
\end{fancydefine}
\begin











\end{document}